%%%%%%%%%%%%%%%%%%%%%%%%%%%%%%%%%
%                               %
%  FDO DJ Opas                  %
%  Akseli Lukkarila             %
%  2025.04.29                   %
%                               %
%%%%%%%%%%%%%%%%%%%%%%%%%%%%%%%%%

\documentclass[12pt, a4paper, oneside]{article}

\usepackage[
    top=2.4cm,
    bottom=1.6cm,
    left=2.1cm,
    right=2.1cm,
    heightrounded
]{geometry}

% Layout
\usepackage[onehalfspacing]{setspace}           % Easy row spacing
\usepackage[parfill]{parskip}                   % Blank line starts new paragraph
\usepackage{fancyhdr}                           % Better header and footer
\usepackage{lastpage}                           % Get total page number
\usepackage{ragged2e}                           % Override ragged commands locally
\usepackage{titlesec}                           % Adjust title and heading spacing
\usepackage{verbatim}                           % Comment sections
\usepackage{layout}                             % Debug margins

% Text
\usepackage{csquotes}                           % Nice quotations
\usepackage{gensymb}                            % Generic symbols for both text and math mode
\usepackage{microtype}                          % Text kerning
\usepackage{soul}                               % Provides hyphenatable spacing
\usepackage{textcomp}                           % Text symbols
\usepackage{units}                              % Nice fractions in text
\usepackage{xurl}                               % Break long urls over lines

% PDF setup and metadata
\usepackage{colorprofiles}
\usepackage[breaklinks, pdfa, pdfversion=1.7]{hyperref}

% Math
\usepackage{unicode-math}                       % Unicode mathematics support for XeTeX and LuaTeX
\usepackage{amsmath}                            % Math stuff

% Colors
\usepackage[luatex,rgb,svgnames,hyperref,table]{xcolor}
\usepackage[skins,breakable]{tcolorbox}

% Tables, lists and figures
\usepackage[inline]{enumitem}                   % Extended list structures
\usepackage{booktabs}                           % Better looking tables
\usepackage{float}                              % Place figure here option [H]
\usepackage{graphicx}                           % External graphics
\usepackage{multicol}                           % Multicolumn pages and lists
\usepackage{multirow}                           % table multirow columns
\usepackage{tabularx}                           % Table sizing options
\usepackage{tikz}                               % tikz graphics

% Captions and subfigures
\usepackage[
    hang,
    bf,
    justification=justified,
    format=plain,
    labelfont={bf},
    textfont=normalfont,
    figurewithin=section,
    tablewithin=section,
]{caption}
\usepackage[labelfont=normalfont,textfont=normalfont]{subcaption}

% Spelling
% Alternative to babel for xelatex and lualatex
% https://ctan.org/pkg/polyglossia
\usepackage{polyglossia}
\setdefaultlanguage{finnish}

% Date formatting
\usepackage[dmyyyy]{datetime}
\renewcommand{\dateseparator}{.}
\usepackage[datesep=.,style=ddmmyyyy]{datetime2}

% Fonts
\usepackage{fontspec}

\newfontfamily\blackfont        [Ligatures=TeX]{Roboto-Black}
\newfontfamily\lightfont        [Ligatures=TeX]{Roboto-Light}
\newfontfamily\semiboldfont     [Ligatures=TeX]{Roboto-Medium}

\setmainfont[Ligatures=TeX,Numbers=Lining]{Roboto}[
    BoldFont = *-Bold,
    BoldItalicFont = *-BoldItalic,
    ItalicFont = *-Italic,
    UprightFont = *-Regular,
]

\AtBeginEnvironment{tabular}{
    \setmainfont[Numbers=Monospaced]{Roboto}[
        BoldFont = *-Bold,
        BoldItalicFont = *-BoldItalic,
        ItalicFont = *-Italic,
        UprightFont = *-Regular,
    ]
}
\AtBeginEnvironment{tabularx}{
    \setmainfont[Numbers=Monospaced]{Roboto}[
        BoldFont = *-Bold,
        BoldItalicFont = *-BoldItalic,
        ItalicFont = *-Italic,
        UprightFont = *-Regular,
    ]
}

\setmathfont{TeX Gyre Termes Math}
\setmonofont{Roboto Mono}

% Path for graphics files
\graphicspath{{./}{kuvat/}}

% Footer
\pagestyle{fancy}
\fancyhf{}
\fancyhead[L]{\small \lightfont FDO DJ OPAS}
\fancyhead[C]{\small \lightfont \today}
\fancyhead[R]{\small \lightfont \thepage}

\renewcommand{\headrulewidth}{0.8pt}            % header line width
\renewcommand{\footrulewidth}{0pt}              % footer line width

\setlength{\headheight}{14pt}
\setlength{\headsep}{12pt}
\setlength{\parindent}{0pt}                     % paragraph indentation
\setlength{\parskip}{12pt plus 2pt minus 2pt}   % paragraph lineskip

% Text and float separation, plus and minus allow for stretching when necessary
\setlength{\textfloatsep}{18pt plus 1pt minus 2pt}
\setlength{\intextsep}{12pt plus 1pt minus 2pt}
\setlength{\abovecaptionskip}{6pt plus 2pt minus 2pt}
\setlength{\belowcaptionskip}{6pt plus 2pt minus 2pt}

\setitemize{label={\LARGE\raisebox{-1.5pt}{\textbullet}}}

% Hyperref settings and PDF metadata
\hypersetup{
    breaklinks=true,
    citecolor=black,
    colorlinks=true,
    linkcolor=black,
    pdfauthor={Akseli Lukkarila},
    pdfmenubar=false,
    pdfpagemode=UseNone,
    pdfstartview=FitH,
    pdfsubject={FDO DJ opas},
    pdftitle={FDO DJ opas},
    pdftoolbar=false,
    unicode,
    urlcolor=NavyBlue,
}

%%%%%%%%%%%%%%%%%%%%%%%%%%%%%%%%%%%%%%%%%%%%%%%%%%%%%%%%%%%%%%%%%%%%%%%%%%%%%%
%%%%%%%%%%%%%%%%%%%%%%%%%%%%%%%%%% DOCUMENT %%%%%%%%%%%%%%%%%%%%%%%%%%%%%%%%%%

\begin{document}

\begin{titlepage}
    \newgeometry{top=6cm, bottom=3cm, right=3.5cm, left=3.5cm}
    \centering
    \includegraphics[width=12cm]{FDO-logo-fix}

    \vspace{1cm}

    {\Huge \blackfont DJ OPAS \par \vspace{6mm}}
    {\Large Musiikin soitto-ohjeet FDO:n kilpailuihin}
    \vfill
    \normalsize
    \today
\end{titlepage}

\cleardoublepage

\doublespacing
\addtocontents{toc}{\protect\enlargethispage{5\baselineskip}}
\tableofcontents
\clearpage

\onehalfspacing

\section{Johdanto} \label{sec:johdanto}

Tämä opas on tarkoitettu ohjeeksi Finnish Dance Organizationin (FDO) alaisissa tanssikilpailuissa musiikin soitosta vastaaville henkilöille.
Oppaassa käydään yksityiskohtaisesti läpi kilpailutapahtuman kulku DJ:n näkökulmasta,
sekä mitä tapahtuman musiikin soittamisessa tulee ottaa huomioon niin esitysmusiikin kuin taustamusiikin osalta.
Koska kyseessä on tanssikilpailu, musiikki on erittäin suuressa roolissa tapahtuman onnistumisen ja sujuvuuden kannalta.
FDO:n kilpailuissa soitettava musiikki voidaan jakaa kolmeen eri osa-alueeseen:
\begin{description}[nosep, labelindent=0.75cm]
    \item [Kilpailumusiikki:] valmiiksi miksatut musiiikikatkelmat, joihin kilpailijat esiintyvät.
    \item [Taustamusiikki:] taukojen aikana soitettava musiikki.
    \item [Palkintomusiikki:] palkintojenjaossa soitettava tunnelmannostatusmusiikki.
\end{description}

Tässä oppaassa keskitytään näistä erityisesti kilpailumusiikkiin,
johon liittyy monia FDO:n kilpailuille ominaisia käytäntöjä sekä tanssilajikohtaisia sääntöja.
Taustamusiikki ja palkintojenjako puolestaan menevät hyvin pitkälti kuten muissa vastaavissa kilpailutapahtumissa,
mutta niissäkin on hyvä ottaa huomioon kilpailtava tanssilaji ja siihen liittyvät musiikkigenret.

Tapahtuman DJ on lähtökohtaisesti vastuussa kaikesta tapahtumassa soitettavasta musiikista,
ellei jotain muuta ole erikseen sovittu.
Tämä tarkoittaa kilpailumusiikkien soittamisen lisäksi taukomusiikeista ja palkintojenjaosta huolehtimista.
Aina kun tapahtumassa ei ole kilpailusarjaa tai juontoa käynnissä,
taustalla olisi hyvä soida kilpailuun sopivaa musiikkia tunnelman ylläpitämiseksi.
Ainoa poikkeus tähän käytäntöön on muodostelmien lavaharjoitus,
joka on yleensä merkitty aikatauluun pidemmän tauon,
kuten ruokailun ajaksi.
Muodostelmien lavaharjoituksen aikana ei soiteta musiikkia.

Kilpailutapahtumassa johtajana toimii \textit{kilpailuvalvoja},
joka nimensä mukaisesti huolehtii kilpailun kulusta ja aikataulusta,
sekä valvoo sääntöjen noudattamista.
Sama henkilö toimii tyypillisesti myös tapahtuman kuuluttajana,
ja istuu samassa rivissä tuomarien pöytien kanssa lavan etureunassa kilpailualueen ja katsomon välissä.
Mikäli tapahtuman aikana herää jotain kysyttävää tai jokin on epäselvää,
kilpailuvalvoja on oikea henkilö vastaamaan kysymyksiin kilpailunjärjestäjän ohella.

Street-kisoissa voidaan kilpailla myös niin sanotuissa battle -sarjoissa,
mutta battlejen soittamista ei käsitellä tässä oppaassa niiden erityisluonteen vuoksi.

DJ-oppaan sisältö on järjestetty seuraavanlaisesti:
Osioissa~\ref{sec:kilpailijan-musiikki} ja~\ref{sec:dj-musiikki} käydään läpi kaksi erilaista esitysmusiikkityyppiä,
jotka ovat kilpailijan oma musiikki ja järjestäjän musiikki.
Osiossa~\ref{sec:taustamusiikki} käsitellään lyhyesti taustamusiikkia ja palkintojenjaon kulkua.
Osiosta~\ref{sec:yhteenveto} löytyy yhteenveto ja tiivistetty muistilista oleellisimmista asioista ja säännöksistä.

\section{Kilpailijan musiikki} \label{sec:kilpailijan-musiikki}

Tietyissä lajeissa ja sarjoissa jokaisella esityksellä on oma valmis esitysmusiikkinsa,
jonka kilpailija on etukäteen toimittanut.
Kilpailussa jokainen tanssija tai ryhmä esiintyy yksi kerrallaan omaan musiikkiinsa.
Musiikkitiedostot ladataan FDO:n nettisivulle ilmoittautumisen yhteydessä,
mistä tiedostot toimitetaan FDO:n toimesta kilpailun DJ:lle keskitetysti kilpailusarjoittain,
yleensä noin viikkoa ennen kilpailutapahtumaa.
\textit{Performing~Arts} -kilpailuissa on käytössä enimmäkseen vain kilpailijan musiikkia,
lukuunottamatta \textit{improvisaatio} -sarjaa, jossa käytetään järjestäjän musiikkia.
\textit{Latino~Show} -kilpailuissa suurimmassa osassa sarjoista on myös käytössä kilpailijan musiikki,
poikkeuksena \textit{Latin Style} -kilpasarjat.
\textit{Disco~Dance}- ja \textit{Street} -lajeissa on puolestaan ainoastaan muutama sarja,
kuten muodostelmat ja katutanssi-showcase,
joissa käytetään kilpailijan omaa musiikkia,
muuten käytössä on järjestäjän musiikki.

Jokaiselle kilpailusarjalle on säännöissä määritetty oma maksimi- ja minimikesto esityksen pituudelle.
Esitysmusiikissa voi olla välissä hiljaisia kohtia,
mutta musiikissa täytyy sääntöjen mukaan olla alku- ja loppuäänet,
jotka merkitsevät esityksen alkua ja loppua.
Etenkin nykytanssiteoksissa saattaa olla pitkiäkin pätkiä hiljaisuutta välissä,
joten älä säikähdä jos musiikki hiljenee yhtäkkiä kesken esityksen.
Usein äänitiedostojen alussa ja lopussa on jonkin verran hiljaisuutta,
mutta musiikki on kuitenkin hyvä soittaa aina aivan alusta asti,
ja erityisen pitkän loppuhiljaisuuden tapauksessa lopettaa kappaleen toisto kesken vasta kun tanssija tai ryhmä on poistunut lavalta.

Mikäli käytät musiikin soittamiseen jotain DJ-ohjelmaa,
kannattaa tiedostot analysoida etukäteen,
ja hyödyntää monista ohjelmista löytyvää automaattista äänenvoimakkuuksien tasausta.
DJ-laitteiden ja ohjelmien kanssa muista varmistaa,
että tempo-liuku on nollassa eikä mitään efektejä ole jäänyt päälle esitysmusiikkeja soitettaessa.
Esimerkiksi CDJ-soittimia tai vastaavia mediasoittimia käytettäessä on suositeltavaa laittaa \textit{tempo~reset} -lukitus päälle esitysmusiikkien soiton ajaksi,
joka estää muutokset kappaleen tempoon.

\subsection{Kilpailujärjestys} \label{subsec:kilpailujarjestys}

Kilpailijoiden musiikkitiedostot toimitetaan jaoteltuina kansioihin niin,
että jokaisen kilpailusarjan musiikit ovat omassa kansiossaan.
Tiedostot on nimetty kilpailunumeron mukaan,
jonka perusteella ne kohdistetaan oikeaan esitykseen.
Jokaisella esityksellä on oma uniikki kilpailunumero,
ja eri sarjojen kilpailunumerot eroavat selkeästi toisistaan.

DJ:lle toimitetaan kilpailun aikana jokaista kilpailukierrosta varten starttilista ennen kyseisen kierroksen alkua,
josta käy ilmi kilpailijoiden esiintymisjärjestys kilpailunumeron mukaan listattuna.
Esitykset ovat aina järjestetty kilpailunumeron mukaan joko nousevasti tai laskevasti.
Esiintymisjärjestys käännetään ympäri aina jokaisen karsintakierroksen jälkeen,
mutta ei kuitenkaan silloin kun ketään ei ole karsittu edelliseltä kierrokselta.
Tälläiset kierrokset on yleensä nimetty \textit{esittelykierros} -nimellä.
Kuvassa~\ref{fig:kilpailunumerot} esitetään esimerkki kilpailunumeroista ja esiintymisjärjestyksestä kolmen kierroksen kilpailusarjassa.

\begin{figure}[ht]
    \centering
    \renewcommand{\arraystretch}{1.0}
    \begin{tabular}{@{\hspace{0.5cm}}c@{\hspace{0.5cm}}}
        \toprule
        Neljännesfinaali        \\ \midrule
        100                     \\
        101                     \\
        102                     \\
        103                     \\
        104                     \\
        105                     \\
        $\vdots$   \vspace{1mm} \\
        120                     \\
        121                     \\ \bottomrule
    \end{tabular}
    \hspace{0.25cm}
    {\huge $\to$}
    \hspace{0.25cm}
    \begin{tabular}{@{\hspace{0.75cm}}c@{\hspace{0.75cm}}}
        \toprule
        Semifinaali             \\ \midrule
        120                     \\
        115                     \\
        114                     \\
        $\vdots$   \vspace{1mm} \\
        108                     \\
        105                     \\
        102                     \\ \bottomrule
    \end{tabular}
    \hspace{0.25cm}
    {\huge $\to$}
    \hspace{0.25cm}
    \begin{tabular}{@{\hspace{1cm}}c@{\hspace{1cm}}}
        \toprule
        Finaali \\ \midrule
        102     \\
        105     \\
        108     \\
        110     \\
        112     \\
        115     \\ \bottomrule
    \end{tabular}
    \caption{Esimerkki esitysmusiikkejen järjestyksestä kilpailukierroksilla.}
    \label{fig:kilpailunumerot}
\end{figure}

\subsection{Kilpakierroksen kulku} \label{subsec:kilpakierroksen-kulku}

On suositeltavaa tehdä jokaisesta kilpasarjasta oma soittolista käytettävään musiikkisoitin- tai DJ-ohjelmaan.
Järjestä biisit tiedostonimen mukaan kierroksesta riippuen joko nousevaan tai laskevaan järjestykseen.
Jokainen esitys kuulutetaan kilpailunumerolla sisään,
ja finaalissa mainitaan lisäksi esityksen ja tanssijan tai ryhmän nimi.
Joskus järjestykseen on voinut tulla äkillisiä muutoksia,
jos esimerkiksi joku tanssija tai ryhmä ei olekkaan vielä valmiina,
tai on joutunut jäämään pois esimerkiksi loukkaantumisen takia.
Tämän vuoksi on tärkeää kuunnella kuuluttajan kertoma kilpailunumero jokaisen esityksen alkuun,
ja näin varmistua seuraavan esityksen numerosta.

Kun kilpanumero on kuulutettu lavalle,
odota rauhassa,
että tanssija(t) pääsevät aloituspaikoilleen ja ovat valmiina aloittamaan ennen musiikin käynnistämistä.
Etenkin isojen muodostelmien lavalletulo ja asettautuminen kestää jonkin aikaa,
joten varmista,
että kaikki tanssijat ovat pysähtyneet lopullisesti paikoilleen ennen musiikin aloittamista.
Lisäksi yleisön hurraus saattaa peittää hiljaisesti alkavan musiikin alun,
joten yleensä on myös parempi odottaa,
että yleisön kannatushuudot ovat laantuneet ennen musiikin aloitusta.
Tarkkaile vielä hetki musiikin käynnistämisen jälkeen,
että esitysmusiikki oli tosiaan oikea.

Jos sarjassa on useampi kilpailukierros,
saat uuden starttilistan ennen saman sarjan seuraavan kierroksen alkua.
Poista nyt aikaisemmasta soittolistasta kaikki ne esitysmusiikit,
jotka eivät päässeet jatkoon,
mutta älä kuitenkaan varmuuden vuoksi poista vielä itse tiedostoja tietokoneelta.
Näin toimimalla oikeat musiikit ovat aina valmiina soittolistassa,
mistä ne on helppo soittaa oikeassa järjestyksessä yksi kerrallaan.

\subsection{Äänenvoimakkuus} \label{subsec:aanenvoimakkuus}

Tyypillisesti esitysmusiikkien välinen äänenvoimakkuus vaihtelee esityksestä toiseen niin teknisistä kuin sisällöllisistä syistä.
Tapahtumassa DJ:n tehtäviin kuuluu oleellisena osana pyrkiä huolehtimaan,
että äänenvoimakkuus pysyy sopivana ja tasaisena kaikkien esitysten kesken.
Toki yksittäisen esitysmusiikin sisällä äänenvoimakkuus voi vaihdella,
mutta yleissääntönä voidaan ajatella,
että yhden kilpailukierroksen aikana jokaisen esitysmusiikin lujimmat osuudet tulisi soida suurin piirtein samalla äänekkyydellä.

Reagoi siis muutoksiin äänenvoimakkuudessa biisien välillä:
Jos jokin biisi soi hiljaisella niin laita sitä lujemmalle,
ja vastaavasti jos jokin biisi tai biisin osuus tulee todella lujaa,
niin laita sitä hiljaisemmalle.
Tällä pyritään varmistamaan ensinäkin,
että mikään esitysmusiikki ei soisi liian suurella äänenvoimakkuudella.
Kilpailutilanteessa on tärkeää pyrkiä tarjoamaan jokaiselle kilpailijalle yhdenvertaiset esiintymispuitteet,
ja näin ollen äänenvoimakkuuksia on hyvä tasoittaa myös reiluuden nimissä.
Tilanne,
jossa joidenkin esityksien musiikki on soinut selkeästi paljon hiljaisemmalla kuin toisten,
ei ole toivottava,
vaikka aina ei pysty täysin tietämään onko jokin osuus kilpailijoiden omasta musiikista tarkoituksella selkeästi hiljaisemmalla,
vai vain teknisen osaamisen puutetta musiikin editoimisessa.

\subsection{Varamusiikki} \label{subsec:varamusiikki}

FDO:n kilpailusääntöjen mukaan kilpailijoilla tulee olla varaversio omasta musiikistaan mukana kilpailupaikalla sen varalta,
että ilmoittautumisen yhteydessä ladattu musiikkitiedosto ei toimi.
Varaversion tulee olla muistitikulla tai laitteella,
josta sen voi soittaa,
kuten älypuhelimella.
Vastaavasti DJ:llä on täytyy olla tekninen valmius varamusiikin soittoon.
On siis suositeltavaa liittää 3,5 millin plugi-kaapeli valmiiksi mikseriin tälläisten tilanteiden varalta.
Mikäli esitysmusiikin toistossa ilmenee ongelmia,
tyypillisesti kyseinen esitys siirretään kierroksen viimeiseksi,
ja kilpailua jatketaan seuraavalla esityksellä sillä välin kun varamusiikki saadaan noudettua ja toistovalmiiksi.

\section{Järjestäjän musiikki} \label{sec:dj-musiikki}

Toinen kilpailuissa käytettävä musiikkityyppi on niin sanottu \textit{järjestäjän musiikki}.
Tämä tarkoittaa,
että kilpailijat tanssivat oman musiikkinsa sijaan DJ:n valitsemaan yhteiseen musiikkiin,
jota kilpailijat eivät etukäteen tiedä.
Musiikkikappaleista soitetaan valmiiksi tehtyjä katkelmia,
joiden pituudet riippuvat kilpailusarjasta ja lajista.
Järjestäjän musiikkia käytetään paljon etenkin \textit{Street}- ja \textit{Disco~Dance} -lajeissa.

Ennen kilpailua täytyy kasata riittävä määrä sopivia kilpabiisejä,
sekä editoida näistä valmiit katkelmat.
Musiikin kasaaminen ja editoiminen voi sopimuksen mukaan olla täysin DJ:n vastuulla,
tai vaihtoehtoisesti DJ:lle voidaan toimittaa joko valmiiksi editoidut katkelmat suoraan,
tai lista sopivista biiseistä,
joista DJ itse editoi katkelmat.
Valmiita musiikkeja tai biisilistoja on mahdollista saada kilpailun järjestävältä tanssikoululta sekä FDO:n lajijaostolta.

Katkelmien tulee olla sanoituksiltaan kilpailuun sopivia,
ja erityisesti lasten ja junioreiden sarjoissa soitettava musiikki tulee olla ikäryhmille sopivaa.
Mikäli DJ vastaa musiikin koostamisesta tai katkelmien tekemistä,
täytyy valmiit biisikatkelmat hyväksyttää FDO:n lajijaostolla (Street-, Disco- tai Latino-show -jaosto kilpailusta riippuen).
Näin varmistutaan,
että biisivalinnat ja katkelmat ovat kilpailuun sopivia.
Musiikin editoimista käsitellään tarkemmin osiossa~\ref{subsec:editoiminen}.

Musiikkikatkelmat on toivottavaa jakaa sarjojen kesken niin,
että \textit{"parhaat"} biisit tulevat finaalikierroksiin,
jotta niin tanssijoiden kuin yleisön energia ja tunnelma ovat silloin parhaimmillaan.
Lapsille ja junioreille on hyvä soittaa hieman nopeampaa ja selkeärytmistä musiikkia,
kun taas aikuisille voi laittaa haastavampia biisejä.
Osiossa~\ref{subsec:rajat} on listattu kootusti kaikkien järjestäjän musiikkia käyttävien sarjojen aika- ja temporajat.

\subsection{Kilpakierroksen kulku} \label{subsec:kilpakierroksen-kulku2}

Järjestäjän musiikkia käyttävissä sarjoissa kilpailukierroksen kulku poikkeaa merkittävästi kilpailijan oman musiikin sarjoista.
Jokaisella kierroksella käytetään kolmea eri musiikkikappaletta,
jotka vuorottelevat kunnes kaikki sarjan kilpailijat ovat tanssineet.
Soolo- ja duosarjoissa on lisäksi alkussa ja loppussa yhteinen tanssikierros neljännellä musiikkikatkelmalla.
Karsintakierroksilla tanssitaan tyypillisesti usean kilpailijan ryhmissä,
mutta finaaleissa jokainen kilpailija tai ryhmä tanssii yksi kerrallaan.

Kun kilpailijoita on paljon,
voidaan kierros jakaa useaan erään (\textit{"heat"}),
jotta kaikki tanssijat mahtuvat yhdessä lavalle: esimerkiksi 48 osanottajan karsintakierroksella voidaan tanssia koko kierros kahdessa osassa niin,
että ensin otetaan lavalle ensimmäiset 24 kilpailijaa,
jotka tanssivat oman yhteisen kierroksen kaikki 24 kilpailijaa lavalla.
Tämän jälkeen loput 24 kilpailijaa tulevat tanssimaan heidän yhteisen kierroksen,
jonka jälkeen ensimmäinen ryhmä tulee tanssimaan pienemmissä ryhmissä oman kierroksensa vuorottelevilla biiseillä,
ja niin eteenpäin.

Samoja biisejä voi käyttää kohtuudella eri kilpailusarjoissa pitkin päivää,
mutta ei kuitenkaan saman sarjan sisällä eri kierroksilla.
Kun musiikkikatkelmat ovat valmiina,
ne on hyvä jakaa valmiiksi kilpailukierroskohtaisiin soittolistoihin ennen tapahtumaa.
Näin varmistetaan,
että jokaiselle kierrokselle tulee sopivat musiikit ja oikeanmittaiset katkelmat,
eikä yhden sarjan sisällä tule samaa biisiä useaan kertaan päivän aikana.
Tapahtumassa edetään nopeasti sarjasta toiseen,
joten kesken kisapäivän on harvemmin aikaa valita lennosta seuraavan kierroksen biisejä.
Erityisesti street-lajeissa on tärkeää,
että yhden kierroksen vuorottelevat kappaleet ovat samantyylisiä ja niissä on sama energia ja tempo.

Kilpailuvalvoja kuuluttaa aina kilpailijat lavalle niin yhteisille kierroksille kuin omalle vuorolle.
Seuraa kuuluttajan ohjeita tanssijoille niin tiedät missä mennään.
Jos kaikki tanssijat ovat lopettaneet ennen katkelmat loppua,
voi sen feidata alas hieman etuajassa.

\subsection{Soolot ja duot} \label{subsec:soolot-ja-duot}

Jokaisella kilpailukierroksella tanssitaan yhteinen kierros,
jossa kaikki kilpailijat tanssivat yhtä aikaa,
sekä toinen henkilökohtainen veto per soolo tai duo.
Karsintakierroksilla pienemmissä ryhmissä muutama kilpailija tai duo kerrallaan,
ja finaaleissa aina yksi kilpailija kerrallaan.
Tätä varten tarvitaan yhteensä neljä eri biisiä jokaiselle kierrokselle:
yksi yhteismusiikki ja kolme vuorottelevaa kappaletta.

\textit{Disco}- ja \textit{Street}-lajeissa sekä \textit{Performing Arts improvisaatio} -sarjassa yhteinen kierros tanssitaan sekä kilpailukierroksen alkuun että loppuun,
eli näissä sarjoissa jokainen kilpailija tai duo tanssii yhteensä kolme kertaa jokaisella kilpailukierroksella.
Molemmilla yhteisillä kierroksilla alussa ja lopussa käytetään samaa musiikkikatkelmaa yhden kilpailukierroksen aikana.
Poikkeuksena \textit{Disco Slow} sooloissa,
jossa alun yhteinen kierros jää pois,
eli tanssitaan vain vuorottelevat kierrokset ja yhteinen loputus.

\textit{Latin Style} -sarjoissa yhteinen kierros tanssivaan vain alussa,
ja finaaleissa tanssitaan ainoastaan yhteinen kierros.

Kilpailusarja alkaa sillä,
että kaikki sarjan kilpailijat otetaan lavalle.
Juontaja kutsuu aina tanssijat esiintymään kilpailunumeroilla.
Odota, että tanssijat ovat levittäytyneet tanssilattialle valmiiksi ennen musiikin käynnistämistä.

Jos kilpailukierros on jaettu useampaan osaan eli \textit{heat}:iin,
jokaiselle \textit{heat}:lle soitetaan samat neljä kappaletta.
Näin tapahtuu yleensä karsintakierroksilla kun osallistujia on niin paljon,
että kaikki eivät mahdu tanssilattialle yhtä aikaa.

\textit{Disco}- ja \textit{Street} -lajeissa jokaisen kierroksen kulku on siis seuraavanlainen:

\begin{itemize}[nosep]
    \item Yhteinen kierros: \textit{biisi D}
    \item Vuorottelevat biisit:
        \begin{itemize}
            \item Kilpailija 1 - \textit{biisi A}
            \item Kilpailija 2 - \textit{biisi B}
            \item Kilpailija 3 - \textit{biisi C}
            \item Kilpailija 4 - \textit{biisi A}
            \item Kilpailija 5 - \textit{biisi B}
            \item Kilpailija 6 - \textit{biisi C}
            \item Kilpailija 7 - \textit{biisi A}
            \item jne\ldots
        \end{itemize}
    \item Yhteinen kierros: \textit{biisi D}
\end{itemize}

\textit{Latin Style}- ja \textit{Performing Arts improvisaatio} -sarjoissa konsepti on sama muutamilla eroavaisuuksilla.

\subsubsection{Disco}

Sooloissa ja duoissa käytetään enimmäkseen yhden minuutin katkelmia muutamaa poikkeusta lukuunottamatta.

Ensimmäinen yhteinen kierros kestää vain 30 sekuntia.

\textit{Disco Slow soolo} -sarjassa alkukarsinnassa ja semifinaalissa yhteiskierros on yhden minuutin mittainen,
mutta finaalissa kilpailijat tanssivat ensin yksi tanssija kerrallaan 45 sekunnin ajan,
ja lopuksi kaikki finalistit tanssivat yhtä aikaa yhden minuutin ajan.

\subsubsection{Street}

\textit{Street} -sarjoissa kaikki musiikkikatkelmat ovat aina yhden minuutin mittaisia sooloissa ja duoissa.

\subsubsection{Performing Arts improvisaatio}

\textit{Improvisaatio} -sarjassa tanssijat kilpailevat soolona.

DJ saa valmiit soittolistat sarjaan Performing Arts -jaostolta tai FDO:n toimistolta.
Yhteen kilpailukierrokseen sisältyy alun yhteinen kierros,
yksitellen tai pienemmissä ryhmissä tapahtuva kierros,
ja sen jälkeen taas yhteinen lopetuskierros.
Jokainen kierros sisältää 1 minuutin kestoisen katkelman,
jotka vaihtuvat joka kerta,
eli jokainen katkelma soitetaan vain kerran tässä sarjassa.
Koko improvisaatio -sarjan aikana yksi sama biisi saa soida ainoastaan yhden kerran.
Sama biisi ei siis saa esimerkiksi soida yhden kerran semifinaali-kierroksella ja toisen kerran finaalikierroksella,
vaan jokaisella kierroksella tulee olla käytössä täysin uudet biisit.
Myös yhteisillä kierroksilla biisien tulee siis vaihtua joka kerta,
eli yhteisillä kierroksilla jokaiselle kierrokselle ja heatille tulee olla eri biisi.
Esimerkiksi kuuden kilpailijan finaalikierrokselle tarvitaan näin ollen kahdeksan eri katkelmaa,
joita ei ole vielä soitettu saman kilpasarjan aikaisemmilla kierroksilla:
alun yhteiskierros,
yksittäiset kierrokset,
ja uusi katkelma vielä lopun yhteiselle kierrokselle.
Tempoja ei tule tasata eri katkelmien kesken tässä sarjassa.

\subsubsection{Latin Style}

Latin Stylessä tanssijat tanssivat kaikilla kilpailukierroksilla kaksi kertaa.
Karsitankierroksilla kaikki kilpailijat tanssivat yhteisen esittelykierroksen yhden minuutin verran,
yleensä useammassa pienemmässä erässä (\textit{heat}).
Tämän jälkeen kilpailijat tanssivat vielä pienemmissä erissä 1,5 minuuttia vuorottelevilla

Finaaleissa tanssitaan vain kerran kaikki yhdessä 1,5 minuuttia.

\subsection{Pienryhmät} \label{subsec:pienryhmat}

Pienryhmillä katkelman pituus on kaksi minuuttia,
paitsi \textit{Disco~Dance Rising Star} ja \textit{Disco~Dance lapset} -pienryhmissä katkelman pituus on yksi ja puoli minuuttia.
Pienryhmäsarjoissa ei tanssita ollenkaan yhteisiä kierroksia,
eli jokainen ryhmä käy tanssimassa yhden kerran yhden kilpakierroksen aikana.
Musiikki toimii samalla periaatteella kuin sooloissa ja duoissa: kolmen biisin kierto jokaisella kilpailukierroksella.
Karsintakierroksilla tanssitaan yleensä kaksi ryhmää kerrallaan,
ja finaalissa yksitellen.
Kilpailijat otetaan lavalla yksitellen.
Odota, että ryhmät ehtivät aloituspaikoilleen ennen musiikin käynnistämistä.
Eli pienryhmillä kierroksen kulku on seuraavanlainen:

\begin{itemize}[nosep]
    \item Vuorottelevat biisit:
        \begin{itemize}
            \item Kilpailija 1 - \textit{biisi A}
            \item Kilpailija 2 - \textit{biisi B}
            \item Kilpailija 3 - \textit{biisi C}
            \item Kilpailija 4 - \textit{biisi A}
            \item Kilpailija 5 - \textit{biisi B}
            \item Kilpailija 6 - \textit{biisi C}
            \item jne\ldots
        \end{itemize}
\end{itemize}

\subsection{Temporajat} \label{subsec:temporajat}

Järjestäjän musiikkia käyttävissä sarjoissa on kilpailumusiikille asetettu laji- ja sarjakohtaiset temporajat (iskua minuutissa).
Temporajat eivät koske kilpailijoiden omaa musiikkia.
Helpoin tapa laskea biisin tempo on analysoida se DJ-ohjelmassa.
Esimerkiksi yleisimmät DJ-ohjelmat \href{https://serato.com/dj}{Serato DJ},
\href{https://rekordbox.com/}{Rekordbox},
ja \href{https://www.virtualdj.com/}{Virtual DJ}
ovat kaikki ilmaisia käyttää \textit{"kirjasto"} -tilassa musiikkikirjaston analysoimiseen ja järjestämiseen ilman kytkettyjä DJ-laitteita.

Temporajat riippuvat sekä ikäluokasta että kilpasarjasta.
Sopivan nopeuden lisäksi kappaleiden tulee tietysti olla muutenkin tanssilajille sopivia.
Seuraavissa osioissa käydään läpi järjestäjän musiikki -sarjojen temporajat lajikohtaisesti.

\subsubsection{Disco Dance}

\textit{Disco~Dance Champion} ja \textit{Disco~Dance Pre-Champion} -sarjoissa musiikin tempo on aina 140 iskua minuutissa.
Tämä tekee katkelmien editoimisesta helppoa,
eikä kilpailukierroksien aikana tarvitse erikseen varmistaa,
että tempot ovat vuorottelevissa katkelmissa samat.
\textit{Disco~Dance Rising Star} -sarjoissa tempo on 136 iskua minuutissa sooloissa ja duossa,
ja 132 iskua minuutissa pienryhmissä.
\textit{Disco Slow} -sarjoissa ei ole varsinaista temporajaa musiikille.

\subsubsection{Latin Style}

\textit{Latin~Style} -sarjoissa kilpaillaan useassa eri lajissa,
joilla jokaisella on oma erityinen musiikkityyli.
Sooloissa kilpaillaan viidessä lajissa: Samba, Cha-Cha, Rumba, Paso Doble, ja Jive.
Duoissa kilpaillaan kahdessa lajissa: Samba ja Cha-Cha.
Kaikissa sarjoissa käytetään järjestäjän musiikkia: \medskip

\begin{table}[ht]
    \centering
    \setlength\tabcolsep{8mm}
    \renewcommand{\arraystretch}{1.5}
    \begin{tabular}{@{\hspace{0.25cm}}ll@{\hspace{0.25cm}}}
        \toprule
        Sarja                  & Tempo                                         \\ \midrule
        Latin~Style Samba      & 100--104 iskua minuutissa (50--52 tahtia/min) \\
        Latin~Style Cha-Cha    & 120--128 iskua minuutissa (30--32 tahtia/min) \\
        Latin~Style Rumba      &  96--108 iskua minuutissa (24--27 tahtia/min) \\
        Latin~Style Paso Doble & 120--124 iskua minuutissa (60--62 tahtia/min) \\
        Latin~Style Jive       & 168--176 iskua minuutissa (42--44 tahtia/min) \\ \bottomrule
    \end{tabular}
    \label{tab:latin-bpm}
\end{table}

\subsubsection{Street}

Street-lajeissa temporajat ovat varsin leveät,
ja niitä kannattaa hyödyntää etenkin usean karsintakierroksen sarjoissa sekä yhteisen katkelman biisivalinnoissa.
Pääpaino tempoissa kannattaa kuitenkin pitää asteikon alemman puolikkaan sisällä,
eli alle 110 iskun minuutissa.

Street-lajeissa kolme vuorottelevaa biisiä tulee olla samantyylisiä yhden kilpailukierroksen sisällä.
Lisäksi ne tulee soittaa samalla tempolla.
Jos esimerkiksi kolme muuten samantyylistä ja toisiinsa hyvin sopivaa biisiä ovat tempoiltaan 100,
104 ja 110 iskua minuutissa,
niin kaikki kolme tulisi soittaa esimerkiksi tempolla 105 iskua minuutissa.
Yhteiskierroksen biisi voi puolestaan olla erilainen ja eri tempolla.
Biisien nopeutuksessa kannattaa kuitenkin pitäytyä kohtuuden rajoissa,
ja nyrkkisääntönä nostaa kappaleiden tempoa maksimissaan kahdeksan prosentin verran.
Alunperin 90 iskua minuutissa olevasta biisistä ei siis kannata väkisin yrittää tehdä duo- tai pienryhmäbiisiä.
Alla olevassa taulukossa esitettään kaikkien Street-sarjojen temporajat: \medskip

\begin{table}[ht]
    \centering
    \setlength\tabcolsep{8mm}
    \renewcommand{\arraystretch}{1.5}
    \begin{tabular}{@{\hspace{0.25cm}}lr@{\hspace{0.25cm}}}
        \toprule
        Sarja                               & Tempo (BPM) \\ \midrule
        Hip~Hop soolo, lapset               & 95--120     \\
        Hip~Hop soolo, juniorit ja aikuiset & 70--116     \\
        Hip~Hop duo                         & 104--120    \\
        Hip~Hop pienryhmä                   & 104--120    \\ \bottomrule
    \end{tabular}
    \label{tab:street-bpm}
\end{table}

\subsection{Musiikin editoiminen} \label{subsec:editoiminen}

Aloita katkelma aina \textit{ykkösellä}, eli tahdin ensimmäisellä iskulla.
Vältä äänenvoimakkuuden häivytyksen (\textit{fade-in}) käyttämistä katkelman alussa,
sillä tällöin tanssijoille on vaikeampaa saada laskuista ja rytmistä kiinni.
Tämä pätee erityisesti \textit{Street}-kisojen musiikkiin.
Katkelman loppuun puolestaan tulee laittaa häivytys (\textit{fade-out}) valmiiksi.
Sopiva loppuhäivytyksen pituus on yhdestä kahteen kasiin (2-4 tahtia),
tai noin viisi sekuntia.

Katkelmaan on hyvä laittaa alkuun 1--2 kasia introa,
joka on mielellään biitillistä osiota eikä pelkkää melodiaa.
Introon voi ottaa esimerkiksi laulutonta osuutta kappaleen alusta tai lopusta,
josta sitten hyppätään suoraan ensimmäiseen säkeeseen.
Lopeta varsinainen tanssiosuus aina täydelle kasille,
vaikka aika menisikin hieman yli.
Käytä samoin introssa ja loppuhäivytyksessä vain täysiä kaseja.
Musiikkikatkelman rakenne kokonaisuudessaan esitetään alla: \\

\begin{figure}[ht]
    \centering
    \includegraphics[width=0.9\textwidth]{katkelma}
    \label{fig:katkelma}
\end{figure}

Intro- ja lopetus-osuuden kanssa katkelmasta tulee siis hieman pidempi kuin nimellinen pituus.
Esimerkiksi yhden minuutin katkelmasta tulee alku- ja loppuosuuksien kanssa noin 1min10s pitkä kokonaisuudessaan.

Katkelmissa on ehdottoman tärkeää välttää kaikenlaisia suvantokohtia ja nostatuksia.
Nyrkkisääntönä yli yhden tahdin mittaiset biittittömät osuudet tulee leikata pois välistä.
Muutenkin kannattaa pyrkiä valikoimaan biisistä aina parhaat palat tiivistettynä katkelmaan,
ja esimerkiksi hyppäämään kertosäkeeseen ennen katkelman loppua.
Mahdolliset kirosanat tulee editoida pois katkelmista.
Samasta biisistä kannattaa tietysti tehdä kaikki soveltuvan mittaiset katkelmat samalla vaivalla: Street-katkelmille yhden ja kahden minuutin versiot (temporajat huomioon ottaen),
ja Discolle kaikki neljä puolesta minuutista kahteen minuuttiin.

\section{Taustamusiikki ja palkintojenjako} \label{sec:taustamusiikki}

Ennen kilpasarjojen alkua ja niiden välissä,
sekä palkintojenjaon alkua odottaessa on hyvä laittaa taustamusiikkia soimaan aina kun juontaja ei ole äänessä.
Kun kilpasarja loppuu niin ole valmiina laittamaan taustamusiikkia soimaan sopivan hetken tullessa.
Pidempien taukojen ja ruokailun ajaksi taustalle voi laittaa soimaan valmiin miksauksen tai soittolistan.
Pääperiaatteena siis taustamusiikin olisi hyvä soida aina kun jotain muuta ei ole meneillään.
Poikkeuksena tähän toimii muodostelmien lavaharjoitus,
jonka aikana taustamusiikkia ei soiteta.
Lavaharjoitus on merkitty kilpailun aikatauluun,
ja se on yleensä joko päivän alussa tai pidemmän tauon aikana.

Kuten kilpailumusiikeissa,
taustamusiikin osalta tulee myös välttää lapsille ja nuorille sopimatonta musiikkia.
Soitettava taustamusiikki on tietysti hyvä sovittaa kilpailtavaan lajiin ja tapahtumaan.
Esimerkiksi Street-kisoissa soita hiphop-musiikkia ja sinne kallellaan olevaa pop-musiikkia.
Disco-kilpailuissa voi soittaa housempaa poppia ja popahtavaa klubimusaa.
Tanssijat myös lämmittelevät tapahtuman taustamusiikkiin etenkin street-kisoissa,
joten aina ennen kilpasarjan alkua kannattaa pyrkiä soittamaan hyvää tanssimusiikkia niin tanssijatkin pääsevät paremmin tunnelmaan sisään ennen lavalletuloa.

\subsection{Palkintojenjako} \label{subsec:palkintojenjako}

Kilpailupäivät päättyvät aina palkintojenjakoon,
missä julkistetaan jokaisen finaalin tulokset ja jaetaan palkinnot sinä päivänä kisatuissa sarjoissa.
Palkinnot jaetaan sarja kerrallaan,
yleensä siinä järjestyksessä kun sarjat on kilpailtu.
Palkintojenjaossa otetaan aina kaikki yhden sarjan kilpailijat aluksi lavalle,
jonka jälkeen tulokset kerrotaan käänteisessä järjestyksessä,
eli viimeiselle sijalle tullut ensimmäisenä ja siitä eteenpäin kohti voittajaa.
Tässä vaiheessa taustalla voi pitää hiljaisella taustamusiikkia pyörimässä.
Erityisesti instrumentaali-versiot biiseistä toimivat hyvin tässä.

Kun voittaja julkistetaan,
odota että juontaja saa kerrottua voittajan tiedot loppuun ja laita sitten soimaan voittobiisi.
Voittobiisi voi olla DJ:n valitsema tai se voidaan nimetä kilpailunjärjestäjän toimesta.
Voittobiisin ei tarvitse olla sama joka sarjassa.
Voittajan julkistamisen jälkeen otetaan yleensä yhteiskuva sarjan kilpailijoiden kanssa,
joten anna voittobiisin soida taustalla kunnes ollaan valmiita siirtymään seuraavaan sarjaan,
tyypillisesti noin minuutin tai pari.
Perinteisiä voittobiisejä on esimerkiksi:
\begin{itemize}[nosep]
    \item Kool \& The Gang -- Celebration
    \item Queen -- We Are The Champions
    \item Irene Cara -- What a Feeling
\end{itemize}

Street-kisoissa hiphop-klassikot toimivat hyvin voittobiiseinä,
kuten esimerkiksi:
\begin{itemize}[nosep]
    \item Will Smith -- Getting Jiggy Wit It
    \item Kanye West -- Touch The Sky
    \item Naughty By Nature -- Hip Hop Hooray
\end{itemize}

\section{Yhteenveto} \label{sec:yhteenveto}

Jos kisapäivän aikana herää mitään kysymyksiä tai jokin on epäselvää,
mene juttelemaan kilpailuvalvojan kanssa.
Tätä on se henkilö,
joka toimii tapahtuman kuuluttajana.
Mikäli tämän oppaan pohjalta jäi jokin asia mietityttämään,
lisätietoja voi kysyä tapahtuman järjestäjältä,
tai perehtyä FDO:n kilpailusääntöihin,
jotka löytyvät FDO:n \href{https://fdo.fi/kilpailut/}{nettisivuilta}.

Tapahtumissa on läsnä lapsia ja nuoria,
joten muista välttää niin kilpailu- kuin taustamusiikeissa avoimen seksuaalista tai paljon kirosanoja sisältävää lyriikkaa.

Seuraavissa ala-osioissa on esitetty tiivistettynä oleellisimmat asiat sekä lajikohtaiset aika- ja temporajat.

\subsection{Kilpailijan musiikki} \label{subsec:kilpailijan-musiikki}

\begin{itemize}[itemsep=-2pt]
    \item Jokaisella esityksella on oma valmis esitysmusiikki.
    \item Kilpailijat esiintyvät yksi kerrallaan.
    \item Esitysmusiikit toimitaan DJ:lle kootusti ennen kilpailua.
    \item Odota, että tanssijat ovat pysähtyneet paikoilleen ennen musiikin aloittamista.
    \item Yritä pitää äänenvoimakkuus tasaisena eri esitysmusiikkien kesken.
    \item Jos soitat DJ-laitteilla tai ohjelmalla, varmista, että tempoliuku (\textit{pitch}) on nollassa eikä mitään efektejä ole päällä.
\end{itemize}

\subsection{Järjestäjän musiikki} \label{subsec:jarjestajan-musiikki}

\begin{itemize}[itemsep=-2pt]
    \item Valmiiksi editoidut lyhyet katkelmat, joihin kaikki kilpailijat tanssivat.
    \item Ei samaa biisiä usealla eri kilpailukierroksella saman sarjan sisällä.
    \item Vuorottelevilla biiseillä tulee olla suurin piirtein sama tempo (sync on kätevä tässä).
    \item Muista huomioida kilpasarjojen temporajat musiikille.
    \item Yritä pitää äänenvoimakkuus tasaisena kaikkien katkelmien kesken.
\end{itemize}

\subsubsection{Performing Arts improvisaatio}

\begin{itemize}[itemsep=-2pt]
    \item Ei samaa biisiä millään kierroksella koko kilpailun aikana.
    \item Biisien tempojen tulee vaihdella laidasta laitaan.
\end{itemize}

\subsection{Tempo- ja aikarajat} \label{subsec:rajat}

Alla olevassa taulukossa esitetään kootusti järjestäjän musiikki -sarjojen musiikkikatkelmien kestot kaikilla osuuksilla,
sekä temporajat (\textit{BPM} eli iskua minuutissa).
Nämä tiedot löytyvät myös kilpailusäännöistä,
ja mikäli näiden kahden dokumentin kesken on eroja,
niin voimassa olevat kilpailusäännöt ovat se mitä noudatetaan.

\begin{table}[ht]
    \centering
    \setlength\tabcolsep{9pt}
    \small
    \renewcommand{\arraystretch}{1.5}
    \begin{tabular}{@{\hspace{0.25cm}}lcccr@{\hspace{0.25cm}}}
        \toprule
        \textbf{Sarja}                                  & \textbf{Yhteisalku} & \textbf{Vuorottelevat} & \textbf{Yhteislopetus} & \textbf{Tempo} \\ \midrule
        Disco Dance soolo                               & 30 s                & 1 min                  & 1 min                  & 140            \\
        Disco Dance duo                                 & 30 s                & 1 min                  & 1 min                  & 140            \\
        Disco Dance pienryhmä                           & ---                 & 1 min 30 s             & ---                    & 140            \\
        Disco Dance Rising Star soolo                   & 30 s                & 1 min                  & 1 min                  & 136            \\
        Disco Dance Rising Star duo                     & 30 s                & 1 min                  & 1 min                  & 136            \\
        Disco Dance Rising Star pienryhmä               & ---                 & 1 min 30 s             & ---                    & 132            \\
        Disco Freestyle soolo                           & 30 s                & 1 min                  & 1 min                  & 140            \\
        Disco Slow soolo                                & ---                 & 1 min                  & 1 min                  & ---            \\ \midrule
        Hip Hop soolo, lapset                           & 1 min               & 1 min                  & 1 min                  & 95--120        \\
        Hip Hop soolo, juniorit ja aikuiset             & 1 min               & 1 min                  & 1 min                  & 70--116        \\
        Hip Hop duo                                     & 1 min               & 1 min                  & 1 min                  & 104--120       \\
        Hip Hop pienryhmä                               & ---                 & 2 min                  & ---                    & 104--120       \\ \midrule
        Latin Style Samba,                              & 1 min               & 1 min 30 s             & ---                    & 100--104       \\
        Latin Style Cha-Cha,                            & 1 min               & 1 min 30 s             & ---                    & 120--128       \\
        Latin Style Rumba,                              & 1 min               & 1 min 30 s             & ---                    &  96--108       \\
        Latin Style Paso Doble,                         & 1 min               & 1 min 30 s             & ---                    & 120--124       \\
        Latin Style Jive                                & 1 min               & 1 min 30 s             & ---                    & 168--176       \\ \bottomrule
    \end{tabular}
    \label{tab:bpm-rajat}
\end{table}

% Uncomment to debug margins
%\clearpage
%\centering
%\layout*

\end{document}
